\chapter{Conclusion}

This thesis investigated whether schema-driven cross-document extraction into a
structured case representation --- a \textit{FactSheet} --- improves the factual accuracy
of an agentic RAG system in multi-document legal settings, and whether it does so without
incurring proportional increases in query-time cost. The motivation was grounded in a
fundamental limitation of chunk-based retrieval: when the answer to a legal query is
distributed across multiple documents, no amount of semantic similarity tuning can make
an isolated-chunk retriever assemble a coherent synthesis. The proposed approach addresses
this by extracting typed entities --- parties, events, claims, damages, deadlines --- from
all case documents at ingestion time and maintaining them as a persistent, relational
representation that supplements retrieval for case-level queries.

\section{Answers to Research Questions}

\paragraph{Main RQ.}
Schema-driven extraction into a structured case representation does improve factual
accuracy: the custom agent achieves statistically significant improvements in correctness
($p = 8.5 \times 10^{-3}$) and relevancy ($p = 3.4 \times 10^{-4}$) relative to the
RAG-only baseline. Token consumption and inference latency are statistically equivalent
between the two architectures, confirming that the quality gain is not purchased at a
proportional resource cost. The main research question is answered affirmatively.

\paragraph{RQ1.}
The performance advantage is largest on the more complex, fully authentic dataset
(TOSL-2024), where queries require synthesising evidence across a larger and more
heterogeneous document set. This dataset-level interaction supports the conclusion that
structured extraction specifically mitigates the cross-document retrieval limitation
rather than providing a diffuse general quality improvement. On the simpler dataset, where
a larger share of queries can be answered from a single source, the gap is smaller ---
consistent with the theoretical expectation that the factsheet adds the most value
precisely where naive retrieval is most likely to fail.

\paragraph{RQ2.}
Session-level trajectories show that the custom agent maintains stable quality with
narrow variance across all sessions as document volume grows. The baseline, by contrast,
exhibits wider fluctuations throughout case progression. The structured representation
scales gracefully with document volume because new information is absorbed incrementally
into the factsheet at each session boundary rather than added to an ever-larger
unstructured retrieval pool. RQ2 is answered affirmatively.

\paragraph{RQ3.}
The model-level breakdown reveals that the performance gap between the custom and
baseline agents is largest for smaller, less capable models and narrows to near-zero
for the most advanced model tested. Structured, pre-digested context reduces the
reasoning burden placed on the underlying model, enabling smaller models to approximate
the performance of larger ones on the RAG-only baseline. This constitutes a practically
significant cost-quality substitution: a cheaper model paired with the custom architecture
can match a more expensive model operating on unstructured retrieval alone. RQ3 is
answered affirmatively, with the caveat that the effect is model-dependent and diminishes
as base model capability increases.

\section{Contributions}

This thesis makes three contributions:

\begin{enumerate}
    \item \textbf{System design.} A schema-driven initialization and incremental update
    pipeline for structured legal case state management, implemented as a LangGraph
    \texttt{StateGraph} with fan-out parallel processing, PostgreSQL relational storage,
    and configurable context compression.

    \item \textbf{Evaluation framework.} A multi-model evaluation setup comparing a
    structured-state agent against a RAG-only baseline across two legal datasets,
    multiple session stages, and five repeated runs per configuration, using a combination
    of semantic quality metrics (DeepEval correctness, relevancy, completeness), surface
    similarity metrics (BERTScore), consistency metrics, and resource consumption measures.

    \item \textbf{Empirical evidence.} Quantitative evidence that structured extraction
    improves factual accuracy in multi-document legal settings without proportional cost
    increase, with the effect being strongest on complex cases, stable across session
    progression, and most pronounced for smaller models.
\end{enumerate}

\section{Limitations and Future Work}

The evaluation rests on two datasets, which limits the generalisability of the findings.
A larger benchmark spanning more cases and document volumes would provide stronger
evidence for the scaling behaviour observed here, and would better distinguish the
effect of document volume from other case-specific characteristics.

The evaluation queries were deliberately restricted to factual, document-grounded
questions to enable objective grading. This scope is appropriate for establishing a
controlled baseline but does not cover the full range of legal reasoning tasks --- in
particular, normative and predictive questions where the benefit of structured extraction
may differ. Extending the evaluation to cover multi-step legal reasoning, conflicting
evidence reconciliation, and temporal argument construction is a natural next step.

The structured representation in this thesis uses a fixed entity schema (parties, events,
claims, damages, deadlines). For cases with different legal structures, such as regulatory
compliance matters or transactional disputes, the schema may require adaptation. An
interesting direction is schema induction: learning the relevant entity types from the
case documents themselves rather than relying on a predefined legal ontology.

The RQ3 result points toward a practically relevant optimization problem: given a target
quality level and a budget constraint, what is the optimal combination of model size and
extraction investment? The evidence suggests this trade-off exists and is exploitable,
but characterising it precisely requires a more systematic sweep over model and
architecture configurations than was feasible within the scope of this thesis.

\paragraph{Advanced retrieval over structured elements.}
The current implementation queries the factsheet through simple SQL lookups filtered by
date range and significance level. The richness of the extracted representation --- typed
entities with source references, timestamps, significance scores, and relational links ---
supports substantially more sophisticated retrieval strategies. Graph traversal across
entity relationships, cross-referencing events with the documents that substantiate them,
and ranking elements by their relevance to a specific query rather than by static
significance are all unexplored. The factsheet as implemented is a data foundation whose
retrieval layer has not been fully exploited.

\paragraph{Agentic reasoning and task queuing.}
The current agent executes a single reasoning pass per query. Adding a task queue that
allows the agent to decompose complex queries into ordered sub-tasks --- each potentially
using a specialised node, toolset, and system prompt --- would support multi-step legal
reasoning that the current flat architecture cannot perform. This is analogous to
chain-of-thought planning at the agent level: rather than producing an answer in one
shot, the agent could plan, retrieve, verify, and synthesise across multiple coordinated
steps.

\paragraph{Background and asynchronous tooling.}
The pipeline already runs offline, but the agent itself is purely reactive: it only acts
when a query arrives. A natural extension is a set of background tools that run during
idle periods --- monitoring connected document repositories and inboxes for new material,
triggering incremental factsheet updates on arrival, or proactively flagging newly
detected deadlines and conflicts to the practitioner. This would move the system from
reactive assistant to proactive case monitor.

\paragraph{Graph database and GraphRAG.}
The relational PostgreSQL schema captures entity types and source links but does not
natively represent or traverse relationships between entities. Migrating the structured
representation to a graph database would enable direct traversal of entity relationships
--- for example, following a chain from a claim to the events that substantiate it to
the documents that evidence those events. Combined with graph-native retrieval, this
would address the cross-document reasoning limitation more completely than the current
hybrid approach and bring the system closer to a full domain-specific GraphRAG
architecture.
